\chapter{Algebry Heytinga}\label{ch:ncl-heyting}

\section{Implikacja jako residuum kresu dolnego}

W poprzednim rozdziale zbudowaliśmy intuicjonistyczny
rachunek zdań i jego semantykę Kripkego
(Twierdzenie~\ref{thm:ncl-ipc-strong-completeness}).
Teraz zbadamy struktury algebraiczne, w których
koniunkcja jest kresem dolnym, alternatywa kresem
górnym, a implikacja jest określona przez relację
sprzężenia porządkowego. Zamiast przypisywać formule
jedną z dwóch wartości klasycznych, będziemy
przypisywać jej element pewnej kraty.

Pojęcia porządku, kresów i kraty wprowadzono w
Definicjach~\ref{def:partial-order},
\ref{def:bounds-sup-inf} oraz \ref{def:lattice}
w logice i teorii mnogości. Korzystamy z praw
działań kratowych z Twierdzenia~\ref{thm:lattice-laws}.
Krata ograniczona ma element najmniejszy $0$
i największy $1$. Dopuszczamy przypadek $0=1$,
czyli kratę jednoelementową.

\begin{definicja}[Algebra Heytinga]
\label{def:ncl-heyting}
Algebrą Heytinga nazywamy strukturę
$H=(H,\wedge,\vee,\to,0,1)$, w której
$(H,\wedge,\vee,0,1)$ jest kratą ograniczoną,
a dla wszystkich $a,b,c\in H$ zachodzi
\begin{equation}
\label{eq:ncl-heyting-adjunction}
 a\wedge c\leq b
 \iff c\leq a\to b.
\end{equation}
Porządek jest porządkiem kraty, tj.
$a\leq b$ wtedy i tylko wtedy, gdy $a\wedge b=a$.
Element $a\to b$ nazywamy względnym
pseudodopełnieniem $a$ względem $b$.
\end{definicja}

Warunek~\eqref{eq:ncl-heyting-adjunction} mówi,
że $a\to b$ jest największą informacją $c$,
która w połączeniu z $a$ wystarcza do uzyskania $b$.
Nie jest to dopełnienie $a$ w sensie klasycznym.
Istnienie takiego największego elementu dla każdej
pary $a,b$ stanowi dodatkowy warunek na kratę.

\begin{lemat}[Jedyność implikacji]
\label{lem:ncl-heyting-unique}
Na ustalonej kracie ograniczonej istnieje co
najwyżej jedna operacja spełniająca
\eqref{eq:ncl-heyting-adjunction}. Istnieje ona
dokładnie wtedy, gdy dla każdej pary $a,b$
zbiór
\[
 D(a,b)=\{c\in H:a\wedge c\leq b\}
\]
ma element największy; wówczas $a\to b=\max D(a,b)$.
\end{lemat}
\begin{proof}
Jeżeli operacja istnieje, to warunek sprzężenia
daje $D(a,b)=\{c:c\leq a\to b\}$, którego
elementem największym jest $a\to b$.
Jedyność wynika z antysymetrii porządku.
Odwrotnie, jeśli $d=\max D(a,b)$, to dla
$c\leq d$ z monotoniczności kresu dolnego
wynika $a\wedge c\leq a\wedge d\leq b$.
Jeżeli $a\wedge c\leq b$, to $c\in D(a,b)$,
więc $c\leq d$. Definiując $a\to b=d$,
otrzymujemy żądane sprzężenie.
\end{proof}

\begin{twierdzenie}[Rozdzielność kraty Heytinga]
\label{thm:ncl-heyting-distributivity}
Każda algebra Heytinga jest kratą rozdzielną.
Ogólniej, jeśli $S\subseteq H$ ma kres górny,
to istnieje kres górny zbioru
$\{a\wedge s:s\in S\}$ i zachodzi
\[
 a\wedge\bigvee S=\bigvee_{s\in S}(a\wedge s).
\]
\end{twierdzenie}
\begin{proof}
Niech $u=\bigvee S$. Z monotoniczności
$a\wedge s\leq a\wedge u$ dla każdego $s\in S$.
Jeśli $b$ jest dowolnym ograniczeniem górnym
zbioru $\{a\wedge s:s\in S\}$, to z residuacji
$s\leq a\to b$ dla każdego $s\in S$.
Stąd $u\leq a\to b$ i ponownie
$a\wedge u\leq b$. To dowodzi wzoru, także
dla $S=\varnothing$, gdy $u=0$.

Dla zbioru dwuelementowego otrzymujemy rozdzielność
kresu dolnego względem kresu górnego.
Druga rozdzielność wynika z pierwszej i praw
kratowych:
\[
\begin{aligned}
 (a\vee b)\wedge(a\vee c)
 &=((a\vee b)\wedge a)\vee((a\vee b)\wedge c)\\
 &=a\vee(a\wedge c)\vee(b\wedge c)\\
 &=a\vee(b\wedge c).
\end{aligned}
\]
W ostatnim kroku użyliśmy prawa pochłaniania.
\end{proof}

\begin{twierdzenie}[Skończone kraty rozdzielne]
\label{thm:ncl-heyting-finite-distributive}
Każda skończona niepusta krata rozdzielna ma
jednoznaczną strukturę algebry Heytinga, określoną
przez
\[
 a\to b=\bigvee\{c:a\wedge c\leq b\}.
\]
\end{twierdzenie}
\begin{proof}
Skończona niepusta krata jest ograniczona:
kres dolny wszystkich elementów daje $0$,
a kres górny daje $1$.
Zbiór $D(a,b)$ jest niepusty, ponieważ zawiera $0$.
Jest skończony, więc ma kres górny $d$.
Rozdzielność, zastosowana indukcyjnie do
skończonej alternatywy, daje
\[
 a\wedge d=\bigvee_{c\in D(a,b)}(a\wedge c)\leq b.
\]
Zatem $d\in D(a,b)$ i jest jego największym
elementem. Stosujemy
Lemat~\ref{lem:ncl-heyting-unique}.
\end{proof}

Skończoność można zastąpić odpowiednim warunkiem
rozdzielności względem dowolnych kresów górnych.

\begin{definicja}[Rama w sensie teorii krat]
\label{def:ncl-heyting-frame}
Ramą nazywamy kratę zupełną $L$, w której dla
dowolnych $a\in L$ i $S\subseteq L$ zachodzi
\[
 a\wedge\bigvee S=\bigvee_{s\in S}(a\wedge s).
\]
Pojęcie to nie oznacza ramy Kripkego
z Definicji~\ref{def:ncl-ipc-kripke}.
\end{definicja}

\begin{wniosek}[Zupełne algebry Heytinga]
\label{cor:ncl-heyting-frames}
Krata zupełna ma strukturę algebry Heytinga
wtedy i tylko wtedy, gdy jest ramą.
Implikację w ramie określa wzór
\[
 a\to b=\bigvee D(a,b).
\]
\end{wniosek}
\begin{proof}
Kierunek prosty wynika z
Twierdzenia~\ref{thm:ncl-heyting-distributivity}.
W ramie kres górny $d=\bigvee D(a,b)$ istnieje
z zupełności, a prawo rozdzielności daje
$a\wedge d=\bigvee_{c\in D(a,b)}(a\wedge c)\leq b$.
Zatem $d$ jest elementem największym $D(a,b)$
i stosujemy Lemat~\ref{lem:ncl-heyting-unique}.
\end{proof}

\section{Podstawowe przykłady i kontrprzykłady}

\begin{przyklad}[Łańcuchy Heytinga]
\label{ex:ncl-heyting-chain}
Każdy ograniczony zbiór liniowo uporządkowany
z działaniami $\wedge=\min$ i $\vee=\max$
jest algebrą Heytinga, jeśli przyjmiemy
\[
 a\to b=
 \begin{cases}
  1,&a\leq b,\\
  b,&a>b.
 \end{cases}
\]
Istotnie, gdy $a\leq b$, każdy $c$ spełnia
$\min(a,c)\leq b$. Gdy $a>b$, nierówność ta
jest równoważna $c\leq b$. W obu przypadkach
wzór wynika z Lematu~\ref{lem:ncl-heyting-unique}.
Na $[0,1]$ otrzymujemy residuum normy trójkątnej
minimum, znane z pierwszego rozdziału.
\end{przyklad}

\begin{przyklad}[Kraty zbiorów domkniętych ku górze]
\label{ex:ncl-heyting-upsets}
Dla zbioru częściowo uporządkowanego $(P,\leq)$
niech $\operatorname{Up}(P)$ będzie rodziną
zbiorów domkniętych ku górze, uporządkowaną przez
inkluzję. Dowolne sumy i przecięcia takich zbiorów
są nadal domknięte ku górze, więc tworzą kratę
zupełną z $0=\varnothing$ oraz $1=P$.
Rozdzielność przecięcia względem dowolnej sumy
wynika z rachunku zbiorów, zatem jest to rama.
Jej implikacja ma postać
\[
 A\to B=\{x\in P:\ \forall y\geq x\
                       (y\in A\Longrightarrow y\in B)\}.
\]
Aby to sprawdzić, oznaczmy prawą stronę przez $D$.
Jeżeli $x\in D$ i $z\geq x$, to każdy $y\geq z$
jest także większy od $x$, więc $z\in D$.
Dla $x\in A\cap D$ przyjmujemy $y=x$ i otrzymujemy
$x\in B$, skąd $A\cap D\subseteq B$.
Jeżeli $C\in\operatorname{Up}(P)$ oraz
$A\cap C\subseteq B$, to dla $x\in C$ i $y\geq x$
mamy $y\in C$. Jeżeli ponadto $y\in A$,
to $y\in B$. Zatem $x\in D$ i $C\subseteq D$.
Stosujemy Lemat~\ref{lem:ncl-heyting-unique}.
\end{przyklad}
Wzór na implikację jest dokładnie klauzulą Kripkego
z Definicji~\ref{def:ncl-ipc-forcing}.
Jeżeli $P=\{r,s\}$ i $r<s$, to mamy trzy elementy
$\varnothing,\{s\},P$, tworzące łańcuch.
Poniższy diagram zestawia ramę Kripkego z jej
algebrą zbiorów domkniętych ku górze.
\begin{center}
\begin{tikzpicture}[>=Stealth]
\node[draw,rounded corners] (r) at (0,0) {$r$};
\node[draw,rounded corners] (s) at (0,1.8) {$s$};
\draw[->,thick] (r)--(s);
\node at (2,0.9) {$\longmapsto$};
\node[draw,rounded corners] (z) at (4,0) {$\varnothing$};
\node[draw,rounded corners,fill=blue!8] (m) at (4,0.9) {$\{s\}$};
\node[draw,rounded corners] (o) at (4,1.8) {$P$};
\draw[thick] (z)--(m)--(o);
\node at (0,-0.65) {Rama Kripkego};
\node at (4,-0.65) {Algebra Heytinga};
\end{tikzpicture}
\end{center}

\begin{przyklad}[Krata rozdzielna bez implikacji Heytinga]
\label{ex:ncl-heyting-nonrelative}
Niech $X=\mathbb N$ oraz
\[
 L=\{F\subseteq X:F\text{ skończony}\}\cup\{X\}.
\]
Jest to krata ograniczona rozdzielna względem
przecięcia i sumy skończonej: te działania
zachowują przynależność do $L$, a prawa
rozdzielności dziedziczymy z kraty potęgowej.
Niech $a=\{1\}$ i $b=\varnothing$.
Zbiór $D(a,b)$ zawiera dokładnie skończone
podzbiory $X\setminus\{1\}$.
Nie ma elementu największego, gdyż każdy
taki podzbiór można powiększyć o kolejny
element $X\setminus\{1\}$.
Z Lematu~\ref{lem:ncl-heyting-unique} wynika,
że na $L$ nie istnieje implikacja Heytinga.
\end{przyklad}
Przykład pokazuje, że sama rozdzielność kraty
ograniczonej nie wystarcza w przypadku nieskończonym.

\section{Rachunek implikacji i pseudodopełnień}

\begin{twierdzenie}[Podstawowe tożsamości implikacji]
\label{thm:ncl-heyting-implication-laws}
Dla dowolnych $a,b,c\in H$ zachodzą
\[
\begin{gathered}
 a\to b=1\iff a\leq b,\qquad
 1\to a=a,\qquad 0\to a=1,\qquad a\to1=1,\\
 a\wedge(a\to b)=a\wedge b,\\
 a\to(b\wedge c)=(a\to b)\wedge(a\to c),\\
 (a\vee b)\to c=(a\to c)\wedge(b\to c),\\
 (a\wedge b)\to c=a\to(b\to c).
\end{gathered}
Implikacja jest antymonotoniczna względem pierwszego
argumentu i monotoniczna względem drugiego.
\end{twierdzenie}
\begin{proof}
Pierwsza równoważność wynika ze sprzężenia dla
$c=1$: $1\leq a\to b$ wtedy i tylko wtedy,
gdy $a\wedge1=a\leq b$.
Równości $0\to a=1$ i $a\to1=1$ są jej
szczególnymi przypadkami. Równość $1\to a=a$
wynika z tego, że $1\wedge c\leq a$ dokładnie
wtedy, gdy $c\leq a$.

Ze sprzężenia mamy $a\wedge(a\to b)\leq b$,
a ten kres dolny jest również nie większy od $a$.
Odwrotnie $a\wedge b\leq b$ daje $b\leq a\to b$,
więc $a\wedge b\leq a\wedge(a\to b)$.

Dla dowolnego $x\in H$ mamy
\[
\begin{aligned}
 x\leq a\to(b\wedge c)
 &\iff a\wedge x\leq b\wedge c\\
 &\iff a\wedge x\leq b\ \text{i}\ a\wedge x\leq c\\
 &\iff x\leq(a\to b)\wedge(a\to c).
\end{aligned}
\]
Równość elementów wynika z antysymetrii:
w równoważności podstawiamy za $x$ kolejno
każdą z obu stron. W analogiczny sposób,
korzystając z rozdzielności, mamy
\[
\begin{aligned}
 x\leq(a\vee b)\to c
 &\iff (a\wedge x)\vee(b\wedge x)\leq c\\
 &\iff x\leq(a\to c)\wedge(b\to c).
\end{aligned}
\]
Łączność kresu dolnego i dwukrotna residuacja dają
\[
 x\leq(a\wedge b)\to c
 \iff b\wedge(a\wedge x)\leq c
 \iff x\leq a\to(b\to c).
\]

Jeżeli $a\leq a'$, to
$a\wedge(a'\to b)\leq a'\wedge(a'\to b)\leq b$,
skąd $a'\to b\leq a\to b$.
Jeżeli $b\leq b'$, to
$a\wedge(a\to b)\leq b\leq b'$,
skąd $a\to b\leq a\to b'$.
\end{proof}

\begin{definicja}[Pseudodopełnienie]
\label{def:ncl-heyting-negation}
W algebrze Heytinga definiujemy
$\neg a=a\to0$. Jest to największy element $c$,
dla którego $a\wedge c=0$, na mocy
Lematu~\ref{lem:ncl-heyting-unique}.
\end{definicja}

\begin{stwierdzenie}[Prawa negacji Heytinga]
\label{prop:ncl-heyting-negation-laws}
Dla dowolnych $a,b\in H$ zachodzą
\[
\begin{gathered}
 \neg0=1,\quad \neg1=0,\quad a\wedge\neg a=0,\\
 a\leq\neg\neg a,\qquad \neg\neg\neg a=\neg a,\\
 \neg(a\vee b)=\neg a\wedge\neg b,\qquad
 \neg a\vee\neg b\leq\neg(a\wedge b).
\end{gathered}
Negacja jest antymonotoniczna.
\end{stwierdzenie}
\begin{proof}
Pierwszy wiersz oraz antymonotoniczność
wynikają z Twierdzenia~\ref{thm:ncl-heyting-implication-laws}
dla następnika $0$.
Ponieważ $\neg a\wedge a=0$, residuacja daje
$a\leq\neg a\to0=\neg\neg a$.
Antymonotoniczność daje $\neg\neg\neg a\leq\neg a$,
a zastosowanie poprzedniej nierówności do $\neg a$
daje nierówność przeciwną.
Pierwsze prawo de Morgana jest instancją tożsamości
$(a\vee b)\to c=(a\to c)\wedge(b\to c)$
dla $c=0$.
Z $a\wedge b\leq a,b$ wynika przez antymonotoniczność,
że $\neg a,\neg b\leq\neg(a\wedge b)$.
Biorąc kres górny, otrzymujemy ostatnią nierówność.
\end{proof}

\begin{twierdzenie}[Kiedy algebra Heytinga jest boolowska]
\label{thm:ncl-heyting-boolean}
Dla algebry Heytinga $H$ następujące warunki
są równoważne:
\begin{enumerate}
\item $H$ jest algebrą Boole'a, tj. kratą ograniczoną
      rozdzielną, w której każdy element ma dopełnienie.
\item Dla każdego $a\in H$ zachodzi $a\vee\neg a=1$.
\item Dla każdego $a\in H$ zachodzi $\neg\neg a=a$.
\end{enumerate}
W takim przypadku $a\to b=\neg a\vee b$.
\end{twierdzenie}
\begin{proof}
Jeżeli $d$ jest dopełnieniem $a$, to $a\wedge d=0$,
więc $d\leq\neg a$.
Z rozdzielności
$\neg a=\neg a\wedge(a\vee d)
=(\neg a\wedge a)\vee(\neg a\wedge d)
=\neg a\wedge d\leq d$.
Stąd $d=\neg a$ i otrzymujemy warunek 2.

Przy warunku 2 rozdzielność daje
\[
 \neg\neg a=\neg\neg a\wedge(a\vee\neg a)
             =\neg\neg a\wedge a\leq a.
\]
Razem z $a\leq\neg\neg a$ daje to warunek 3.
Jeżeli zachodzi warunek 3, to
\[
 \neg(a\vee\neg a)=\neg a\wedge\neg\neg a=0.
\]
Zatem $a\vee\neg a=\neg\neg(a\vee\neg a)=1$.
W warunku 2 element $\neg a$ jest dopełnieniem
$a$, ponieważ $a\wedge\neg a=0$, a rozdzielność
została dowiedziona w
Twierdzeniu~\ref{thm:ncl-heyting-distributivity}.
Otrzymujemy warunek 1.

Na koniec, przy warunku 2, jeśli $a\wedge c\leq b$,
to
$c=c\wedge(a\vee\neg a)
=(c\wedge a)\vee(c\wedge\neg a)\leq b\vee\neg a$.
Odwrotnie $c\leq\neg a\vee b$ daje
$a\wedge c\leq a\wedge(\neg a\vee b)=a\wedge b\leq b$.
Jedyność residuum kończy dowód.
\end{proof}

\begin{przyklad}[Trójelementowa algebra nieboolowska]
\label{ex:ncl-heyting-three-chain}
W łańcuchu $0<u<1$ z
Przykładu~\ref{ex:ncl-heyting-chain} mamy
$\neg u=0$ i $\neg\neg u=1$.
Zatem $u\vee\neg u=u\neq1$ oraz
$\neg\neg u\neq u$.
Tablica implikacji, z poprzednikiem w wierszu
i następnikiem w kolumnie, ma postać
\[
\begin{array}{c|ccc}
 \to&0&u&1\\\hline
 0&1&1&1\\
 u&0&1&1\\
 1&0&u&1
\end{array}
\]
W szczególności implikacja nie jest operacją
$\neg a\vee b$: dla $a=b=u$ pierwsza daje $1$,
druga $u$. Jest to algebra trójwartościowa Gödlowska,
a nie macierz Kleene'ego z drugiego rozdziału.
\end{przyklad}

\begin{stwierdzenie}[Związek z kratami rezyduowanymi]
\label{prop:ncl-heyting-residuated}
Algebry Heytinga są dokładnie ograniczonymi
przemiennymi całkowitymi kratami rezyduowanymi
(z Definicji~\ref{def:lst-fuzzy-residuated-lattice}),
w których działanie monoidalne jest idempotentne.
\end{stwierdzenie}
\begin{proof}
W algebrze Heytinga działanie $\wedge$ jest łączne,
przemienne i idempotentne, ma element neutralny $1$,
a residuacja jest warunkiem definicyjnym.
Odwrotnie niech $*$ oznacza działanie monoidalne
w takiej kracie rezyduowanej. Z residuacji wynika
monotoniczność: jeżeli $u\leq v$, to
$u\leq v\leq a\to(a*v)$, więc $a*u\leq a*v$.
Przemienność daje monotoniczność w obu argumentach.
Ponieważ $a,b\leq1$, mamy $a*b\leq a*1=a$
oraz $a*b\leq1*b=b$, zatem $a*b\leq a\wedge b$.
Dla $d=a\wedge b$ idempotentność i monotoniczność
dają $d=d*d\leq a*b$.
W konsekwencji $a*b=a\wedge b$ i residuacja
jest dokładnie warunkiem algebry Heytinga.
\end{proof}
Nie każda algebra Heytinga spełnia prelinearność,
wymaganą w algebrach MTL. Przykładowo w algebrze
$\operatorname{Up}(P)$ dla rozgałęzionej ramy
z Przykładu~\ref{ex:ncl-ipc-fork-countermodel}
przyjmijmy $A=\{a\}$ i $B=\{b\}$.
Korzeń $r$ nie należy ani do $A\to B$, ani do $B\to A$,
więc $(A\to B)\vee(B\to A)\neq P$.
